\section{2014级工科数学分析下 B卷}

\subsection{资料来源}
\begin{itemize}
  \item 试题：2014级软件工科数学分析下 B 卷，DOC 正文已抽取；公式对象已按 Word 公式图片逐项恢复。
  \item 公式图片审计：\texttt{tmp/past-exam-equations/2014-B/manifest.csv}。
\end{itemize}

\subsection{试题}
诚信应考，考试作弊将带来严重后果！

华南理工大学本科生期末考试

《工科数学分析》2014--2015 学年第二学期期末考试试卷（B）卷

注意事项：
\begin{enumerate}
  \item 开考前请将密封线内各项信息填写清楚；
  \item 所有答案请直接答在试卷上（或答题纸上）；
  \item 考试形式：闭卷；
  \item 本试卷共 5 个大题，满分 100 分，考试时间 120 分钟。
\end{enumerate}

\subsubsection*{一、填空题（每小题 4 分，共 20 分）}
\begin{enumerate}
  \item 微分方程
  \[
  y''-y=2xe^x
  \]
  的特解形式为 \underline{\hspace{4cm}}。
  \item 设由方程
  \[
  f(x-y,y-z,z-x)=0
  \]
  确定了 $z=z(x,y)$，其中 $f$ 有连续的偏导数，则 $dz=$ \underline{\hspace{3cm}}。
  \item 交换积分的顺序，则
  \[
  \int_0^1dy\int_{y^2}^{\sqrt y}f(x,y)\,dx=
  \]
  \underline{\hspace{3cm}}。
  \item 设 $f(x)$ 是周期为 $2$ 的周期函数，它在区间 $(-1,1]$ 上的表达式为
  \[
  f(x)=
  \begin{cases}
  3,&-1<x\le0,\\
  x^3,&0<x\le1,
  \end{cases}
  \]
  则 $f(x)$ 的 Fourier 级数在 $x=1$ 处收敛于 \underline{\hspace{3cm}}。
  \item 设
  \[
  L:|x|+|y|=1
  \]
  取逆时针方向，则
  \[
  \int_L xy^2\,dx+x^2y\,dy=
  \]
  \underline{\hspace{3cm}}。
\end{enumerate}

\subsubsection*{二、计算题（每小题 10 分，共 40 分）}
\begin{enumerate}
  \item 设
  \[
  z=f\!\left(2x,\frac{y}{x}\right),
  \]
  其中 $f$ 具有二阶连续的偏导数，求 $\dfrac{\partial^2z}{\partial x\partial y}$。
  \item 计算三重积分
  \[
  I=\iiint_\Omega z\,dx\,dy\,dz,
  \]
  其中 $\Omega$ 是由球面 $x^2+y^2+z^2=4$ 与抛物面 $x^2+y^2=3z$ 所围成的区域。
  \item 计算曲线积分
  \[
  I=\int_L\frac{-y\,dx+x\,dy}{x^2+4y^2},
  \]
  其中 $L$ 是从 $A(-\pi,0)$ 沿曲线 $y=\cos\dfrac{x}{2}$ 到 $B(\pi,0)$ 的有向曲线弧段。
  \item 设 $\Sigma$ 是锥面 $z=\sqrt{x^2+y^2}$ 被平面 $z=0$ 及 $z=1$ 所截部分的下侧，计算曲面积分
  \[
  I=\iint_\Sigma 2x\,dy\,dz+3y\,dz\,dx+(z^2-5z)\,dx\,dy.
  \]
\end{enumerate}

\subsubsection*{三、解答题（每小题 10 分，共 20 分）}
\begin{enumerate}
  \item 求幂级数
  \[
  \sum_{n=1}^{\infty}\frac{2n-1}{n!}x^{2n}
  \]
  的和函数。
  \item 设
  \[
  f(x)=\sin x-\int_0^x(x-t)f(t)\,dt,
  \]
  其中 $f(x)$ 为连续函数，求 $f(x)$。
\end{enumerate}

\subsubsection*{四、证明题（本题 10 分）}
证明：级数
\[
\sum_{n=1}^{\infty}\frac{\sin(nx)}{n^2}
\]
在 $(-\infty,+\infty)$ 上一致收敛。

\subsubsection*{五、应用题（本题 10 分）}
求椭球面
\[
x^2+y^2+\frac{z^2}{4}=1
\]
在第一卦限的一点，使该点处的切平面在三个坐标轴上的截距平方和最小。

\subsection{答案与解析}

\subsubsection*{一、填空题}
\begin{enumerate}
  \item 特征方程 $r^2-1=0$，$r=1$ 是一重特征根。右端为一次多项式乘 $e^x$，故特解形式为
  \[
  y^*=xe^x(Ax+B).
  \]
  \item 记 $f_1,f_2,f_3$ 为 $f$ 对三个自变量的偏导数。由
  \[
  f_1(dx-dy)+f_2(dy-dz)+f_3(dz-dx)=0
  \]
  得
  \[
  dz=\frac{(f_3-f_1)\,dx+(f_1-f_2)\,dy}{f_3-f_2}.
  \]
  \item 区域为 $0\le y\le1,\ y^2\le x\le\sqrt y$。等价于 $0\le x\le1,\ x^2\le y\le\sqrt x$，故
  \[
  \int_0^1dx\int_{x^2}^{\sqrt x}f(x,y)\,dy.
  \]
  \item 周期为 $2$，在 $x=1$ 处左右极限分别为 $x^3\to1$ 与周期延拓后的 $3$，故 Fourier 级数收敛于
  \[
  \frac{1+3}{2}=2.
  \]
  \item 由 Green 公式，
  \[
  \int_Lxy^2\,dx+x^2y\,dy=\iint_D(2xy-2xy)\,dA=0.
  \]
\end{enumerate}

\subsubsection*{二、计算题}
\begin{enumerate}
  \item 令 $\xi=2x,\ \eta=y/x$。先有
  \[
  z_y=\frac1x f_\eta.
  \]
  再对 $x$ 求偏导，得
  \[
  z_{xy}=\frac2x f_{\xi\eta}-\frac{y}{x^3}f_{\eta\eta}-\frac1{x^2}f_\eta,
  \]
  其中各偏导在 $\left(2x,\dfrac{y}{x}\right)$ 处取值。
  \item 用柱坐标。两曲面交于 $r^2+z^2=4,\ r^2=3z$，得 $z=1,\ r=\sqrt3$。所围区域为
  \[
  0\le r\le\sqrt3,\qquad \frac{r^2}{3}\le z\le\sqrt{4-r^2},\qquad 0\le\theta\le2\pi .
  \]
  因此
  \[
  I=\int_0^{2\pi}\int_0^{\sqrt3}\int_{r^2/3}^{\sqrt{4-r^2}}z r\,dz\,dr\,d\theta
  =\frac{13\pi}{4}.
  \]
  \item 令 $u=x,\ v=2y$，则
  \[
  \frac{-y\,dx+x\,dy}{x^2+4y^2}
  =\frac12\frac{-v\,du+u\,dv}{u^2+v^2}.
  \]
  沿题设方向从 $A(-\pi,0)$ 到 $B(\pi,0)$ 时，向量 $(u,v)$ 的极角从 $\pi$ 变到 $0$，故
  \[
  I=\frac12(0-\pi)=-\frac{\pi}{2}.
  \]
  \item 下侧锥面 $z=r$ 对应法向取
  \[
  (z_x,z_y,-1)=\left(\frac{x}{r},\frac{y}{r},-1\right).
  \]
  投影区域 $D:x^2+y^2\le1$，且 $z=r$。于是
  \[
  I=\iint_D\left(2x\frac{x}{r}+3y\frac{y}{r}-(r^2-5r)\right)\,dx\,dy.
  \]
  用极坐标积分：
  \[
  I=\int_0^{2\pi}\int_0^1\left(2r\cos^2\theta+3r\sin^2\theta-r^2+5r\right)r\,dr\,d\theta
  =\frac{9\pi}{2}.
  \]
\end{enumerate}

\subsubsection*{三、解答题}
\begin{enumerate}
  \item 令 $t=x^2$。则
  \[
  \sum_{n=1}^{\infty}\frac{2n-1}{n!}x^{2n}
  =2\sum_{n=1}^{\infty}\frac{n t^n}{n!}-\sum_{n=1}^{\infty}\frac{t^n}{n!}
  =2te^t-(e^t-1).
  \]
  因此和函数为
  \[
  (2x^2-1)e^{x^2}+1.
  \]
  \item 对积分方程求两次导数。先有
  \[
  f'(x)=\cos x-\int_0^x f(t)\,dt,
  \]
  再得
  \[
  f''(x)=-\sin x-f(x).
  \]
  即
  \[
  f''+f=-\sin x.
  \]
  由原式得 $f(0)=0,\ f'(0)=1$。通解为
  \[
  f=C_1\cos x+C_2\sin x+\frac12 x\cos x.
  \]
  代入初值，得
  \[
  f(x)=\sin x+\frac12 x\cos x.
  \]
\end{enumerate}

\subsubsection*{四、证明题}
对任意 $x\in(-\infty,+\infty)$，
\[
\left|\frac{\sin(nx)}{n^2}\right|\le\frac1{n^2}.
\]
而 $\sum_{n=1}^{\infty}\dfrac1{n^2}$ 收敛，由 Weierstrass 判别法，原级数在 $(-\infty,+\infty)$ 上一致收敛。

\subsubsection*{五、应用题}
椭球面可写为 $F=x^2+y^2+\dfrac{z^2}{4}-1=0$。在点 $(x_0,y_0,z_0)$ 处切平面为
\[
x_0x+y_0y+\frac{z_0z}{4}=1.
\]
三个截距为
\[
\frac1{x_0},\qquad \frac1{y_0},\qquad \frac4{z_0}.
\]
需在约束 $x_0^2+y_0^2+\dfrac{z_0^2}{4}=1$ 下最小化
\[
\frac1{x_0^2}+\frac1{y_0^2}+\frac{16}{z_0^2}.
\]
令 $u=x_0^2,\ v=y_0^2,\ w=z_0^2/4$，则 $u+v+w=1$，目标为
\[
\frac1u+\frac1v+\frac4w.
\]
Lagrange 乘子法给出 $u=v,\ w=2u$，故 $4u=1$。第一卦限内所求点为
\[
\left(\frac12,\frac12,\sqrt2\right).
\]
